\documentclass[9pt,a4paper]{extarticle}
\usepackage[utf8]{inputenc}
\usepackage[margin=0.45in,top=0.42in,bottom=0.42in]{geometry}
\usepackage{xcolor}
\usepackage{hyperref}
\usepackage{booktabs}
\usepackage{tabularx}
\usepackage{enumitem}
\usepackage{titlesec}
\usepackage{fancyhdr}
\usepackage{mathptmx}

\definecolor{primary}{RGB}{20, 70, 115}
\definecolor{secondary}{RGB}{40, 120, 180}
\definecolor{darkgray}{RGB}{50, 50, 50}

\hypersetup{colorlinks=true, linkcolor=primary, urlcolor=secondary}

\titleformat{\section}{\small\bfseries\color{primary}}{}{0em}{}[\titlerule]
\titlespacing*{\section}{0pt}{4pt}{2pt}

\pagestyle{fancy}
\fancyhf{}
\lhead{\footnotesize\textcolor{darkgray}{\textbf{IoTrix 2.0} $|$ Track A: Embedded IoT System}}
\rhead{\footnotesize\textcolor{darkgray}{Team: \textbf{Neural-Nexus} $|$ \textbf{IntelliCane}}}
\lfoot{\footnotesize\textcolor{darkgray}{\href{https://github.com/Neural-Nexus-IoTrix-2-0/intellicane}{github.com/Neural-Nexus-IoTrix-2-0/intellicane}}}
\rfoot{\footnotesize\textcolor{darkgray}{One-Page Technical Proposal}}
\renewcommand{\headrulewidth}{0.4pt}
\renewcommand{\footrulewidth}{0.4pt}

\setlist{nosep, leftmargin=1.2em}

\begin{document}

\begin{center}
    {\Large \textbf{\color{primary} IntelliCane: One-Page Technical Proposal}} \\[0.2em]
    {\normalsize \textbf{Smart Assistive Cane with Haptic Obstacle Avoidance, Fall Detection, \\ and a Companion App for Blind Users and Caretakers}} \\[0.3em]
    \footnotesize
    \textbf{Team Neural-Nexus} \quad$|$\quad \textbf{Track A: Embedded IoT System} \quad$|$\quad \textbf{IoTrix 2.0 Semi-Final (5 Oct 2026)} \\
    \textbf{Members:} Dulnith Liyanage, Thenul Sahansa, Chamith Chethana, Suneth Vidurasa \quad$|$\quad \textbf{Cost:} LKR 2,800 \\
    \textbf{Repository Link:} \url{https://github.com/Neural-Nexus-IoTrix-2-0/intellicane}
\end{center}
\vspace{-0.6em}
\hrule height 0.8pt
\vspace{0.4em}

\section{1. The Problem: Why We Built IntelliCane}
Navigating everyday environments with a standard white cane is tough. Visually impaired individuals frequently bump into knee-to-head level hazards (like open windows, branches, or table corners) and risk serious injury if they drop their cane or fall. Existing commercial smart canes often cost over \$500 (LKR 150,000+), depend entirely on internet connections, or blast loud beeps that drown out traffic and ambient sound. As undergraduate students, our goal was to build a practical, reliable, and truly affordable (LKR 2,800) smart cane that keeps users safe offline while connecting them to family and caretakers.

\section{2. Our Proposed Solution}
\textbf{IntelliCane} pairs an instant offline hardware reflex on the cane with an accessible Android companion app (\texttt{com.example.intellicane}):
\begin{itemize}
    \item \textbf{Offline Safety Reflex (ESP32-C3):} Runs completely on the cane with a sub-50\,ms loop. An HC-SR04 ultrasonic sensor scans ahead up to 4\,m, driving a handle vibration motor via 200\,Hz PWM that gently starts at 60\,cm and smoothly intensifies as obstacles get closer (down to 10\,cm). An MPU6050 IMU detects sudden drops and falls ($>65^\circ$ tilt or $>2.5\,g$ impact). When a fall happens, the cane stays quiet locally (no loud beeping or buzzing to panic the fallen user) and immediately sends emergency alert packets over BLE.
    \item \textbf{Companion Android App (\texttt{application/}):} Connects over BLE 5.0 (Nordic UART Service) and provides two modes: \textit{Blind User Mode} (clean accessible dashboard, one-tap location sharing, and a collapsible floating terminal for live debugging) and \textit{Caretaker Mode} (real-time presence status, Firebase Firestore fall sync, and an interactive OSMDroid/ESRI map with an instant FAB that centers on the user with 18.5x zoom using the phone's native GPS).
\end{itemize}

\section{3. System Architecture \& Data Flow}
\vspace{-0.2em}
\begin{center}
\footnotesize
\begin{tabular}{|l|c|l|c|l|}
\hline
\textbf{Sensing (Cane)} & $\rightarrow$ & \textbf{Embedded Reflex (ESP32-C3)} & $\rightarrow$ & \textbf{Tactile \& Audio Alerts} \\
HC-SR04 (Trig: 0, Echo: 1 direct) & & 160 MHz RISC-V, Non-blocking scheduler & & 200 Hz LEDC PWM Haptic Motor (GPIO 6) \\
MPU6050 6-Axis IMU (I2C: 4/5) & & Monotonic ramping math \& fall debounce & & Proximity Audio Warning (Silent on Fall) \\
\hline
\multicolumn{5}{|c|}{\textit{Wireless BLE 5.0 Nordic UART Service (RX: 6E400002 / TX: 6E400003)}} \\
\hline
\textbf{Mobile App (Kotlin)} & $\rightarrow$ & \textbf{Location \& Cloud Services} & $\rightarrow$ & \textbf{Caretaker View} \\
User Dashboard \& Caretaker Modes & & Smartphone A-GPS Geolocation & & Real-time OSMDroid/ESRI Map (FAB Centering) \\
Floating Terminal \& BLE Client & & Firebase Firestore Sync \& Presence & & Remote Fall Alert Push \& Presence \\
\hline
\end{tabular}
\end{center}
\vspace{-0.4em}

\section{4. What We've Built \& Verified So Far}
\begin{itemize}
    \item \textbf{Working Hardware Bench Prototype:} Built on an ESP32-C3 SuperMini with HC-SR04 Echo wired directly to GPIO 1, GY-521 MPU6050, 3-pin vibration motor module, and buzzer.
    \item \textbf{Smooth Proportional Haptics:} 200\,Hz PWM vibration scales smoothly from 60\,cm down to 10\,cm, staying completely off past 60\,cm so the user's hand doesn't tire out.
    \item \textbf{Silent Fall Detection over BLE:} Tilting past $65^\circ$ or impact $>2.5\,g$ silences the local motor and buzzer, immediately transmitting \texttt{FALL\_STATE:1} over BLE; standing the cane back up ($<30^\circ$) sends \texttt{FALL\_STATE:0} and resets navigation.
    \item \textbf{Resilient I2C Sensor Driver:} Custom direct-register fallback driver handles clone MPU6050/6500 chips, multiple I2C addresses (\texttt{0x68}/\texttt{0x69}), and recovers automatically from bus lockups.
    \item \textbf{Live Android Mobile App:} Native Kotlin app (\texttt{com.example.intellicane}) with User Dashboard, collapsible floating terminal, Firestore cloud fall sync (\texttt{updateFallStateInFirestore}), and live map with user-centering FAB (\texttt{fabUserLocation}).
\end{itemize}

\section{5. Components \& Budget Breakdown (Phase 1 Bench: LKR 2,800)}
\begin{itemize}
    \item \textbf{Hardware Parts:} ESP32-C3 SuperMini (LKR 1,400), HC-SR04 Ultrasonic (LKR 450), GY-521 MPU6050 IMU (LKR 600), 3-Pin Vibration Motor (LKR 250), Piezo Buzzer \& Jumpers (LKR 100). \textbf{Total: LKR 2,800} (well within our LKR 20,000--35,000 project budget ceiling).
    \item \textbf{Software Stack:} C++17, Arduino-ESP32, Kotlin Android, Firebase Firestore \& RTDB, OSMDroid / ESRI, BLE 5.0 Nordic UART Service.
\end{itemize}

\section{6. Testing What We Built}
\begin{itemize}
    \item \textbf{Host Unit Tests (\texttt{tests/proximity\_feedback\_test.cpp}):} 100\% pass across boundary inputs (negative, NaN, $\infty$), smooth PWM ramping, motor polarity flips, beep transitions, and 32-bit \texttt{millis()} rollover safety.
    \item \textbf{Bench Calibration:} Ultrasonic sensor verified with physical tape measurements from 2\,cm to 250\,cm with $\pm 1.5$\,cm accuracy. Fall trigger verified with a digital protractor ($>65^\circ$ trigger, $<30^\circ$ clear).
    \item \textbf{Wireless BLE Performance:} Stable 10\,Hz telemetry stream to smartphone with $<25$\,ms latency over a 10\,m range. Electronics stayed cool ($<32^\circ$C) during continuous vibration testing.
\end{itemize}

\section{7. Known Limitations \& How We're Handling Them}
\begin{itemize}
    \item \textbf{Ultrasonic Reflections on Soft Clothes:} Sound waves can scatter on angled soft fabrics beyond 2\,m (\textit{Fix:} Adding a compact Time-of-Flight laser sensor for the final build).
    \item \textbf{Breadboard Jumpers:} Prototype uses jumpers that can loosen with heavy motion (\textit{Fix:} Designing a custom 2-layer PCB and 3D-printed handle before the final).
    \item \textbf{Phone Battery Dependency for Cloud:} If the phone dies, cloud syncing pauses (\textit{Safety Reflex:} The cane's core obstacle sensing remains 100\% functional offline).
\end{itemize}

\section{8. Roadmap to the Final (17 October 2026)}
1. \textbf{Custom PCB \& 3D Handle:} Design a soldered PCB and 3D-print an ergonomic handle housing with TPU shock damping.
2. \textbf{Automated Caretaker Alerts:} Wire up Firebase Cloud Functions to dispatch automated SMS and push notifications on confirmed fall events.
3. \textbf{Battery Optimization:} Integrate an 18650 Li-ion battery with TP4056 charge controller and implement ESP32-C3 Light Sleep for all-day runtime.

\end{document}
